%% ALMOST DONE -- edit only where something stated here has actually changed, and then minimally.
%% docs_style.md section 1 is the rule.
\section{The long run}\label{sec:sp:longrun}

The exhaustibility of discoveries gives the model a hard long-run question. Cumulative virgin extraction is bounded by~\eqref{eq:sp:sum:cumulative-N}, so the economy cannot draw on nature forever; what it can do is stretch a finite material budget through the recycling loop, hold material in its two anthropogenic stocks --- the capital stock and the waste stock --- and, if technology permits, keep a closed loop turning forever. With the minimum material requirement read as a hard floor, the question sharpens from ``what does the economy converge to'' to ``does the economy survive'': there is no dematerialized mode to retire into, and holding the floor forever is the only alternative to collapse. This section asks which long-run states are feasible and which are optimal to end up in. Three primitives organise the answers: whether production has a material floor ($\bar R=0$ or $\bar R>0$), whether the yield ceiling is hard or soft, and whether technology is asymptotically stationary or trending. Their combinations span a grid of cases, whose entries we call \emph{cells}; the goal of the section is to attach to each cell the long-run state the economy ends up in. We proceed through a small set of maintained scoping assumptions (Section~\ref{sec:sp:longrun:scope}), a block of accounting facts that hold on every feasible path (Section~\ref{sec:sp:longrun:budget}), a classification of the candidate terminal states available to fill the cells (Section~\ref{sec:sp:longrun:class}), their characterisation with and without the floor (Sections~\ref{sec:sp:longrun:floor}--\ref{sec:sp:longrun:trend}), and the taxonomy table that assembles the full grid (Section~\ref{sec:sp:longrun:taxonomy}). Appendix~\ref{sec:app:workhorse} sets up a workhorse specification with explicit functional forms; we point to it throughout.

\subsection{Scope: maintained long-run assumptions}\label{sec:sp:longrun:scope}

The system of Section~\ref{sec:sp:opt} carries six states, is non-autonomous through the technology trends, is non-convex through the decay function ($\theta'<0$) and the material floor, and switches between corner and interior regimes on several margins. We start by restricting the scope of the long run analysis with some basic assumptions. Beyond ruling out limit cycles, history-dependent states and other complex theories, we will be explicit about what the assumptions rule out and what they do not.

Throughout this section, growth and decay rates are per-period log rates: a variable $Z$ ``grows at rate $g$'' when $\ln Z_{t+1}-\ln Z_t\to g$, and ``decays at rate $\nu$'' when it grows at rate $-\nu$, so that products of trending variables simply add their rates.

\begin{assumption}[technology regimes]\label{ass:lr:tech}
One of the following holds.
\begin{enumerate}
    \item[(S)] \emph{Asymptotically stationary technology.} Every primitive converges to a limit function (uniformly on compact sets): $F\left(\cdot,t\right)\to F^\infty$, $a\left(\cdot,t\right)\to a^\infty$, $C^N\left(\cdot,t\right)\to C^{N\infty}$, $C^D\left(\cdot,t\right)\to C^{D\infty}$, $c^c(t)\to c^{c\infty}$, $c^T\left(\cdot,t\right)\to c^{T\infty}$, and each limit inherits the properties assumed in Section~\ref{sec:sp:setup:primitives}, with the yield ceiling in either case (H) or (S).
    \item[(T)] \emph{Trending technology.} Final-goods technology combines TFP growth and material-augmenting progress at constant rates,
    \begin{align}
        F\left(K^Y,R,P,t\right) &= A_t\,\tilde{\mathcal F}\left(K^Y,B_tR^{e},P\right)\ \text{for }R\ge\bar R,\quad 0\text{ otherwise} \label{eq:sp:lr:trendF} \\
        \text{and }\frac{A_{t+1}}{A_t}&=e^{g_A},\ g_A\ge0, \qquad
        \frac{B_{t+1}}{B_t}=e^{g_B},\ g_B\ge0, \notag
    \end{align}
    with $g_A+g_B>0$, and all remaining primitives --- in particular the recycling technology and the handling share $\mu$ --- asymptotically stationary as in (S).
\end{enumerate}
\end{assumption}

Regime (S) is the reading of $F_t\ge0$, $a_t\ge0$, $C^N_t\le0$ in which improvements eventually peter out: since the trends are monotone, boundedness is all that needs to be added. Regime (T) isolates the two engines of progress that matter for the material question --- making output from less input altogether ($g_A$), and making output from less \emph{physical} input in particular ($g_B$) --- while holding the recycling technology asymptotically fixed, which keeps the circularity margins from becoming a third engine. One feature of~\eqref{eq:sp:lr:trendF} should be marked at once, because much of what follows turns on it: material-augmenting progress multiplies the \emph{effective} input $R^e=R-\bar R$, never the floor. Technology can economise without bound on the material used above $\bar R$; it cannot economise on $\bar R$ itself, which is a statement about matter, not about knowledge.

\begin{assumption}[bounded accumulation and finite values]\label{ass:lr:bounded}
In regime (S), accumulation alone cannot sustain growth: $\lim_{K\to\infty}F^\infty_K\left(K,R,0\right)<\delta$ uniformly for $R$ in compact sets. In both regimes, discounted values are finite along the candidate paths considered: lifetime utility converges, and every costate admits its forward (fundamental) representation. In regime (T), with $g$ denoting the asymptotic growth rate of consumption and $\eta$ the elasticity of marginal utility, this requires $\ln\left(1+\rho\right)>\left(1-\eta\right)g$.
\end{assumption}

\begin{assumption}[convergence]\label{ass:lr:convergence}
Optimal paths settle down. In regime (S), states and controls converge to a rest point of the limit system. In regime (T), every variable converges after detrending: goods-block variables ($C$, $K$, $Y$, cost flows) grow at constant asymptotic rates, and material-block variables (tonnage flows and the stocks $S$, $\bar X_{\max}-X$, $M^K$, $\mathcal W$) decay at constant asymptotic rates or converge.
\end{assumption}

\begin{assumption}[regeneration floor]\label{ass:lr:theta}
Decay is bounded away from zero, $\theta\left(P\right)\ge\theta_{\min}>0$ for all $P$, and valuation never tips: $\rho+\theta\left(P\right)+\theta'\left(P\right)P>0$ on the relevant range of $P$.
\end{assumption}

Four comments on what has been assumed away. First, Assumption~\ref{ass:lr:convergence} excludes limit cycles and more exotic attractors outright; the pollution feedback $\theta'<0$ is exactly the kind of mechanism that can generate them, and we make no claim that they cannot occur in this model --- only that we do not study them. It does \emph{not} exclude multiplicity of rest points or history dependence: when several long-run states are consistent with the conditions below, which one the economy reaches may depend on initial stocks, and we condition on the limit attained. Second, Assumption~\ref{ass:lr:bounded} rules out explosive accumulation in the stationary regime, so that ``sustained consumption'' means a level, not a rate. Third, Assumption~\ref{ass:lr:theta} makes pollution fully reversible: once emissions cease, the stock regenerates completely. Permanent pollution legacies --- $\theta\to0$ at finite $P$ --- are a genuine possibility we set aside; allowing them would add a hysteresis dimension to every entry of the classification below. Fourth, a remark on ranking collapse paths: several cells below have $C\to0$, and if $u$ is unbounded below ($u\left(0\right)=-\infty$, as under CRRA with $\eta\ge1$) then every path in such a cell has $U_0=-\infty$ and the criterion cannot rank them. Whenever we speak of the optimal path \emph{within} a collapse cell --- the timing of the shutdown, the value extracted on the way down --- we implicitly assume $u$ bounded below, or rank by an overtaking criterion; this is a genuine restriction, and the quantitative model must respect it.

\subsection{The material budget and four accounting facts}\label{sec:sp:longrun:budget}

The results of this subsection require no optimality and almost no limit-taking: they are accounting, and they hold on every feasible path in both regimes. Define the \emph{material budget}
\begin{align}
\mathcal B &\equiv \left(1+\Omega^{N,S}\right)\left(S_0+\bar X_{\max}-X_0\right)+\Omega^{D,S}\left(\bar X_{\max}-X_0\right)+M^K_0+\mathcal W_0,
\label{eq:sp:lr:budget}
\end{align}
in tonnes: everything that can ever be inside the economy's circulation or enter its waste stream from outside --- extractable material and the overburden it brings up, the mass displaced by all the exploration the ceiling permits, and the material already held in the two anthropogenic stocks at date zero.

\begin{lemma}[stock bound and throughput cap]\label{lem:lr:stocks}
On every feasible path and in every period,
\begin{align}
M^K_t+\mathcal W_t \;\le\; \mathcal B,
\qquad\text{and hence}\qquad
R_t \;\le\; N_t+\bar a\,\mu\,\mathcal B .
\label{eq:sp:lr:stockbound}
\end{align}
\end{lemma}

The collected ledger~\eqref{eq:sp:setup:ledger-collected} reads $\Delta\left(M^K+\mathcal W\right)_{t+1}=N_t+\mathcal N_t-\left(1-a_t\varpi_t\right)H_t\le N_t+\mathcal N_t$; summing and using the cumulative bounds on $N$ and $D$ from~\eqref{eq:sp:sum:cumulative-N} and $X_{t+1}-X_t=D_t$ gives the first inequality. The second is the throughput cap~\eqref{eq:sp:setup:throughput-cap} with the stock replaced by its bound. \emph{Circular supply is bounded in every period, under both ceiling cases}: however close $a\varpi$ comes to one, recycled material cannot exceed $\bar a\mu$ tonnes per tonne of retained anthropogenic stock, and the stock can never exceed the budget.

\begin{lemma}[cumulated leakage]\label{lem:lr:leakage}
On every feasible path and for every horizon $T$,
\begin{align}
\sum_{t=0}^{T} \left(1-a_t\varpi_t\right)H_t \;\le\; \mathcal B .
\label{eq:sp:lr:leakage}
\end{align}
\end{lemma}

This is the same summation read on the other side: cumulated leakage equals cumulated inflows plus the net rundown of the two stocks, both bounded by the budget. Lemma~\ref{lem:lr:leakage} holds under both ceilings and is the general form from which everything ceiling-specific follows.

\begin{lemma}[finite cumulative throughput under the hard ceiling]\label{lem:lr:throughput}
Under the hard ceiling (H), on every feasible path,
\begin{align}
\sum_{t=0}^\infty H_t \;\le\; \frac{\mathcal B}{1-\bar a},
\qquad
\sum_{t=0}^\infty R_t \;\le\; \left(S_0+\bar X_{\max}-X_0\right)+\frac{\bar a\,\mathcal B}{1-\bar a}
\;\equiv\;\bar{\mathcal R},
\qquad
\sum_{t=0}^\infty W_t \;\le\; \frac{\mathcal B}{1-\bar a}+\mathcal B .
\label{eq:sp:lr:throughput}
\end{align}
If in addition $\bar R>0$, the material era has bounded duration: the number of periods in which production runs is at most
\begin{align}
\#\left\{t:R_t\ge\bar R\right\} \;\le\; \bar{\mathcal R}/\bar R .
\label{eq:sp:lr:duration}
\end{align}
\end{lemma}

The first bound follows from Lemma~\ref{lem:lr:leakage} and $1-a\varpi\ge1-\bar a>0$; the second from $R=N+a\varpi H$; the third from $W=\Delta\mathcal W+H$ and the stock bound; the last from $\sum_t R_t\ge\bar R\cdot\#\left\{R_t\ge\bar R\right\}$. This is the circularity multiplier in cumulated form: recycling lets each tonne of the budget serve production at most $\left(1-\bar a\right)^{-1}$ times, and no feasible policy --- however capital-intensive its recycling, however patient its planner --- can do better. Under the hard ceiling, physical throughput, waste and pollution are therefore \emph{transitional} phenomena on every feasible path: the flows are summable, under Assumption~\ref{ass:lr:convergence} they converge to zero, and with the regeneration floor $P\to0$.

The soft ceiling is another matter, and the next result says exactly what escaping Lemma~\ref{lem:lr:throughput} costs.

\begin{proposition}[growth is a precondition for perpetual circularity]\label{prop:lr:growth}
Suppose a feasible path holds $R_t\ge\bar R>0$ for all $t\ge T_0$. Then:
\begin{enumerate}
\item the yield ceiling is soft ($\bar a=1$) and $\sum^\infty\left(1-a_t\varpi_t\right)<\infty$, so $a\varpi\to1$ and recycling capital per treated tonne diverges, $x\to\infty$;
\item the retained stock covers the floor forever,
\begin{align}
\mathcal W_t \;\ge\; \frac{\bar R-N_t}{\bar a\,\mu}
\qquad\text{with }N_t\to0,
\label{eq:sp:lr:survival}
\end{align}
so that $\liminf_{t\to\infty}\mathcal W_t\ge\bar R/\left(\bar a\,\mu\right)=\bar R/\mu$ under the soft ceiling: the anthropogenic reserve, turning over at share $\mu$ per period, must itself carry the floor;
\item recycling capital, and with it total capital and output, grow without bound. Under Assumption~\ref{ass:lr:bounded} this is infeasible in regime (S): \emph{perpetual circularity requires trending technology with $g>0$.}
\end{enumerate}
\end{proposition}

The proof chains the accounting. Since $\sum_t N_t<\infty$, eventually $N_t<\bar R/2$ and recycled supply must cover the rest, $a\varpi\mu\mathcal W\ge\bar R-N$, which is~\eqref{eq:sp:lr:survival} and bounds $H=\mu\mathcal W$ away from zero. Lemma~\ref{lem:lr:leakage} with $H$ bounded below forces $\sum\left(1-a\varpi\right)<\infty$, which is impossible under a hard ceiling and, under the soft ceiling, forces $a\to1$ and hence $x=K^R/\left(\varpi H\right)\to\infty$ with $\varpi H$ bounded below --- so $K^R\to\infty$. Maintaining an unbounded capital stock requires unboundedly growing replacement output, which Assumption~\ref{ass:lr:bounded} rules out in regime (S). The proposition is the model's central feasibility fact: \emph{a closed material loop is not a low-tech steady state but a capital-intensive frontier that only a growing economy can afford.} Which tails of the yield function make the diverging recycling intensity worth paying for is a value question, not a feasibility one; it is taken up with the taxonomy below and settled quantitatively in the workhorse.

\subsection{Candidate terminal states}\label{sec:sp:longrun:class}

Classify a terminal configuration by what the material block and the goods block do. The material block can be \emph{dead} (all material flows zero from some date, or below the floor so that production has ceased), \emph{vanishing} (flows strictly positive forever but converging to zero), a \emph{perpetual loop} ($N\to0$ with $R$ bounded away from zero, sustained by recycling), or \emph{perpetual extraction} ($N$ bounded away from zero). The goods block can collapse ($C\to0$), settle at a rest point ($C\to C^*>0$), or grow ($g>0$). The accounting eliminates most combinations:

\begin{itemize}
\item \emph{Perpetual extraction} is eliminated outright by the cumulative bound~\eqref{eq:sp:sum:cumulative-N}, in every cell.
\item \emph{A perpetual loop with a stationary or collapsing goods block} is eliminated by Proposition~\ref{prop:lr:growth}: neither a bounded nor a shrinking economy can maintain the diverging recycling capital a closed loop requires.
\item With the hard floor, \emph{a vanishing material block} is possible only when $\bar R=0$: with $\bar R>0$, flows below the floor mean production has already ceased, so the material era ends discretely or not at all.
\end{itemize}

Five named states survive, and they organise the rest of the section. We tag them by letter group --- A for the collapse of a floor economy, B for the floorless benchmark states, C for circular growth. With a floor ($\bar R>0$): \textbf{collapse} (A: the material era ends in finite time by the discrete shutdown of Section~\ref{sec:sp:opt:foc}, and with it all production) and \textbf{perpetual circular growth} (C: the loop closes and rides on trending technology). Without a floor ($\bar R=0$), three benchmark states: \textbf{asymptotic squeeze to a dematerialized rest point} (B1), \textbf{balanced dematerialization} under trending technology (B2), and \textbf{exhaustible-resource collapse} (B3).

\subsection{Long-run states with a material floor}\label{sec:sp:longrun:floor}

With $\bar R>0$ the hard reading of the floor does two things at once: it makes materials essential in level by construction and it pre-empts marginal essentiality as the shutdown occurs at a strictly positive flow, before any Inada divergence at the threshold can be reached along an optimal path. Neither essentiality notion classifies anything here; the classification collapses onto two primitives, the ceiling and the regime.

\begin{proposition}[no productive rest point]\label{prop:lr:norest}
Let $\bar R>0$. At any rest point of the limit system, $D=N=W=H=R=0$, $\mathcal W=0$, $M^K=0$, $P=0$, and $Y=C=0$. Consumption can be positive in the long run only if it grows, there is no stationary survival.
\end{proposition}

The proof runs the rest-point conditions through the stocks one at a time. $\Delta X=0$ forces $D=0$; $\Delta S=0$ then forces $N=0$, so $\mathcal N=0$. With $\Delta M^K=\Delta\mathcal W=0$ the collected ledger gives $\left(1-a\varpi\right)H=0$; under the hard ceiling this forces $H=0$ directly, and under the soft ceiling $a\varpi=1$ requires $x=\infty$, unavailable at any rest point with finite capital, so $H=0$ again. Then $\mathcal W=H/\mu=0$, $W=\Delta\mathcal W+H=0$, $R=N+a\varpi H=0<\bar R$, hence $Y=0$ and, with gross investment non-negative, $C=0$. Finally $\Xi=0$ and the regeneration floor give $P=0$. Under regime (S), Assumption~\ref{ass:lr:convergence} then makes \emph{collapse the only possible long run with a floor}, whatever the ceiling: this is the floor's version of the survival dichotomy, with the dichotomy's benign branch deleted.

\smalltitle{Collapse (A).} In every cell with $\bar R>0$ except one, the terminal state is collapse, and the accounting says how it arrives. Under the hard ceiling (or any path on which cumulative throughput is finite), Lemma~\ref{lem:lr:throughput} bounds the duration of the material era: production can run at or above the floor for at most $\bar{\mathcal R}/\bar R$ periods in total. The era does not fade out --- the floor forbids running material use continuously to zero --- but ends by the discrete shutdown of Section~\ref{sec:sp:opt:foc}: production runs at or above the floor for as long as it runs, and stops permanently once holding the floor is no longer feasible. After the shutdown $Y=0$, consumption is zero, and the capital stock runs off.

Appendix~\ref{sec:app:workhorse:shutdown} makes both halves of this explicit for the workhorse. During the era the material value obeys Hotelling's rule applied to the recycling-multiplied budget, $\Psi_{t+1}/\Psi_t=1+r_{t+1}$, so circularity lowers the level of the price path without touching its slope --- the precise sense in which it is a transition technology here. After the shutdown the economy is a cake-eating problem in its own depreciating capital stock, with consumption declining geometrically at factor $\left[\left(1-\delta\right)/\left(1+\rho\right)\right]^{1/\eta}$ rather than vanishing at once; that closed form is what makes the value of stopping computable, and hence the shutdown date itself.

Two features of the aftermath deserve record. First, \emph{trending technology does not rescue the economy}: $B_t$ multiplies $R-\bar R$ and cannot economise on the floor, so growth changes the level of consumption extracted during the material era --- possibly enormously --- and the optimal timing of the shutdown, but not the verdict. Collapse occurs \emph{despite} growth in the cell (H)$\times$(T). Second, the economy leaves a pollution echo. The stockpile cannot be emptied in finite time ($\mathcal W_t\ge\left(1-\mu\right)^t\mathcal W_0>0$), and after the shutdown there is no output to pay for treatment, so $\varpi=0$ and the draining stock passes to the environment untreated: $\Xi_t=\mu\mathcal W_t$, with the stock running down geometrically at factor $1-\mu$. Collapse is followed by a legacy emission stream --- landfills outlive the economy that filled them --- after which $P\to0$ by the regeneration floor.

\smalltitle{Perpetual circular growth (C).} The single surviving cell is $\bar R>0$, soft ceiling, regime (T). Its configuration follows from the results above: $N\to0$ and $D\to0$ (the resource side dies as always); $R$ converges to a strictly positive level $R_\infty>\bar R$ met entirely by recycled material; the loop closes, $a\varpi\to1$ with $x\to\infty$; the stockpile converges to the level that turns over the loop, $\mathcal W_\infty=R_\infty/\mu$, satisfying the survival bound~\eqref{eq:sp:lr:survival}; leakage $\left(1-a\varpi\right)\mu\mathcal W$ is summable, so $\Xi\to0$ and $P\to0$; the material intensity $\Omega\propto R/Y\to0$; and the goods block grows, in the workhorse's Cobb--Douglas case at
\begin{align}
g &= \frac{g_A+\gamma\,g_B}{1-\beta_K},
\label{eq:sp:lr:c2-growth}
\end{align}
with $\beta_K$ and $\gamma$ the output elasticities of capital and materials --- the balanced rate of an economy whose physical material input is asymptotically \emph{constant} while its effective input $B_tR^e_\infty$ grows at $g_B$. Consumption grows without bound on a fixed diet of matter: the circular economy here is not an efficiency improvement but the entire survival condition, and by Proposition~\ref{prop:lr:growth} it is affordable only because the economy grows.

The construction behind this configuration is carried out for the workhorse forms in Appendix~\ref{sec:app:workhorse:cbgp}. Four of its results belong in the main text, because they are what turns the cell from a description into a characterisation.

\emph{The circulating flow obeys Little's law, and survival is a stock condition.} Summing the collected ledger over all periods defines the \emph{retained material endowment} $\mathcal M_\infty\equiv\lim_t\left(M^K_t+\mathcal W_t\right)\le\mathcal B$ --- the budget less what was never extracted and what leaked on the way. On the path the two anthropogenic stocks are stationary, so each is its flow times its residence time, and
\begin{align}
R_\infty &= \frac{\mathcal M_\infty}{\mathcal T},
&
\mathcal T &\equiv \frac{1}{\mu}+\frac{\sigma_\infty}{\delta},
&&\text{so survival requires}
&
\mathcal M_\infty &> \bar R\,\mathcal T,
\label{eq:sp:lr:little}
\end{align}
with $\mathcal T$ the mean number of periods a tonne spends standing still per pass --- waiting in the stockpile, and locked in durables at the storage share $\sigma_\infty$. Since $\mathcal M_\infty\le\mathcal B$ and $\mathcal T\ge1/\mu$, a necessary condition in primitives alone is $\mu\mathcal B>\bar R$: the entire material budget, turning over at the handling share, must clear the floor. The direction of the durability term is worth marking, because it is counter-intuitive: for a given retained endowment, a \emph{more} durable economy supports a \emph{smaller} sustainable throughput, since a larger share of its matter is not moving.

\emph{The price system is bounded, and waste is priced as a perpetuity.} As $a\varpi\to1$ the effective survival factor $1-\mu\left(1-\alpha\varpi\right)$ of~\eqref{eq:sp:sum:Wstock} tends to one, so a stockpiled tonne is a claim on every future delivery rather than on one, and $p^W=-\Theta_{\mathcal W}\Psi$ with $\Theta_{\mathcal W}=\mu e^{g}/\left(1+r-e^{g}\right)$, finite exactly under Assumption~\ref{ass:lr:bounded}. All prices then grow at rate $g$ while $\Omega\propto R/Y$ decays at $g$, so every wedge $\zeta\Omega^j$ is a bounded constant --- the boundedness the construction needed --- and the whole asymptotic block reduces to a closed-form scalar system. One sign is worth recording: $\zeta<0$. Embodied material carries a \emph{negative} shadow price on a closed loop, because holding matter in durables keeps it out of the recovery loop where it is worth $\Theta_{\mathcal W}\Psi$ per pass. Waste is the last mine, and at the margin durability competes with it.

\emph{The recycling margin is satisfiable, and what limits it is the tail, not the ceiling.} Along the path $\tilde F_K$ is constant while $\mathcal G^K$ grows at $g$, so the capital margin $a'\left(x_t\right)\mathcal G^K_t=\tilde F_K$ has a unique root $x_t$ for all large $t$ and $x_t\to\infty$ automatically; the apparent degeneracy of the exponential tail dissolves once $\Psi$ is read as driven by calendar time rather than by $x$. What is not automatic is that the leakage the path leaves behind is cumulatively finite, as Lemma~\ref{lem:lr:leakage} demands. Appendix~\ref{sec:app:workhorse:cbgp} shows that the per-period leakage rate is the recycling capital share deflated by the elasticity of the yield tail, $1-a\left(x_t\right)\propto s^R_t/\epsilon_a\left(x_t\right)$ with $\epsilon_a\equiv xa'/\left(1-a\right)$, and that summability of this product --- not the shape of the tail as such --- is the criterion. Both parametric families in general use pass it: the exponential tail closes the loop with $x_t$ rising \emph{linearly} in time and a recycling sector whose capital share vanishes, and the power tail closes it for every exponent, because the value of closing the shortfall grows forever while the cost of doing so does not. Tails do exist that fail --- logarithmic ones, where $\epsilon_a\to0$ --- and for those a floor economy collapses exactly as under a hard ceiling. But the speed at which the ceiling is approached is not what decides the cell; the height of the ceiling is, and the taxonomy below is stated accordingly.

\emph{Selection.} Where state~C is feasible and the alternative is collapse, it is selected whenever $u$ ranks a growing stream above one converging to zero, which is the honest sense in which it is the survival route. Where $\bar R=0$ and the alternative is the balanced dematerialization path of Section~\ref{sec:sp:longrun:trend}, the comparison is settled by~\eqref{eq:app:wh:bdp-vs-c}: the closed loop grows strictly faster, at asymptotically vanishing capital cost, so the planner closes the loop whenever the technology allows it. What the floor changes is not the choice but the stakes.

One caveat bounds the whole construction. The feasibility side --- Proposition~\ref{prop:lr:growth}, \eqref{eq:sp:lr:little} and the closure criterion --- is accounting plus a margin that has been solved. Sufficiency is not established: the problem is non-convex through $\theta'<0$, through the floor and through the corner structure, so what Appendix~\ref{sec:app:workhorse:cbgp} delivers is a verified candidate path satisfying every necessary condition, not a uniqueness theorem. That gap is recorded among the open items below and is not closed by anything in this paper.

\subsection{The floorless benchmark: stationary technology}\label{sec:sp:longrun:restpoints}

The case $\bar R=0$ is the natural benchmark: its states isolate what the floor itself is doing, and they are the limits the economy approaches when $\bar R$ is small. With $\bar R=0$ the specification is the standard $Y=\mathcal F\left(K^Y,R,P,t\right)$, dematerialized production exists whenever $\mathcal F\left(K,0,P\right)>0$, and essentiality is a genuine parameter again.

\begin{lemma}[material autarky at rest]\label{lem:lr:autarky}
At any rest point of the limit system, $D=N=W=H=R^R=R=\Xi=0$, $\mathcal W=0$, $\Omega=0$, $M^K=0$ and $P=0$. Consumption at a rest point with capital $K$ is
\begin{align}
C &= F^\infty\left(K,0,0\right)-\delta K .
\label{eq:sp:lr:restC}
\end{align}
\end{lemma}

The proof is that of Proposition~\ref{prop:lr:norest} without the floor's last step, plus the intensity condition~\eqref{eq:sp:intensity} giving $\Omega=0$ (whenever $\mathcal D+\phi^IG>0$) and $\Delta M^K=-\delta M^K=0$ giving $M^K=0$. Every rest point is a point on the dematerialized technology, and the entire materials block --- the budget, the recycling technology, the handling share, the waste parameters --- drops out of the long-run state. Whether a good rest point exists is then a property of the dematerialized technology alone.

\begin{definition}[asymptotic essentiality]\label{def:lr:essentiality}
Let $\bar C\equiv\sup_{K\ge0}\left[F^\infty\left(K,0,0\right)-\delta K\right]$. Materials are \emph{asymptotically inessential in level} if $\bar C>0$, and \emph{asymptotically essential} if $F^\infty\left(K,0,P\right)=0$ for all $K,P$. Materials are \emph{asymptotically essential at the margin} if $F^\infty_R\left(K,R,0\right)\to\infty$ as $R\to0$ for relevant $K$, and \emph{inessential at the margin} if this limit is finite (a finite choke value).
\end{definition}

The level notion decides survival; the marginal notion decides how the material era ends. The two are independent, and the workhorse spans the combinations with one substitution parameter. Note the contrast with Section~\ref{sec:sp:longrun:floor}: both notions matter \emph{only} when $\bar R=0$.

\begin{proposition}[survival dichotomy, stationary technology, $\bar R=0$]\label{prop:lr:dichotomy}
Under Assumptions~\ref{ass:lr:tech}(S) and \ref{ass:lr:bounded}--\ref{ass:lr:theta}, with $\bar R=0$:
\begin{enumerate}
    \item If materials are asymptotically inessential in level, there is a rest point with positive consumption (state B1). The best such rest point satisfies the modified golden rule of the dematerialized technology,
    \begin{align}
    F^\infty_K\left(K^*,0,0\right) &= \rho+\delta,
    &
    C^* &= F^\infty\left(K^*,0,0\right)-\delta K^*>0,
    \label{eq:sp:lr:mgr}
    \end{align}
    with limit prices $r=\rho$, $q=1$, $u'\left(C^*\right)=\lambda^*$, $p^{S*}=p^{X*}=0$,
    \begin{align}
    p^{P*} &= \frac{v'\left(0\right)/\lambda^*-F^\infty_P\left(K^*,0,0\right)}{\rho+\theta\left(0\right)},
    &
    p^{M*} &= \frac{\delta}{\rho+\delta}\,p^{W*},
    &
    \zeta^* &= \sigma^*\frac{\rho}{\rho+\delta}\,p^{W*},
    \label{eq:sp:lr:restprices}
    \end{align}
    where $\sigma^*$ is the limit of the storage share in the form~\eqref{eq:sp:zeta-explicit}, and $p^{W*}=-p^{\mathcal W*}=-\mu h^*/\left(\rho+\mu\right)$ is determined by the rest-point form of the waste margin, \eqref{eq:sp:sum:W-limit}. Because $\Omega=0$, every material wedge $\zeta\Omega^j$ vanishes and the rest-point allocation is that of a standard Ramsey model.
    \item If materials are asymptotically essential in level, every rest point has $C=0$, and under Assumption~\ref{ass:lr:convergence} consumption converges to zero on every optimal path: collapse (state B3). Recycling and the material budget determine the value and duration of the decline, not its occurrence.
\end{enumerate}
\end{proposition}

Part (i) collects the limits of the optimality conditions of Section~\ref{sec:sp:opt:foc}. Two of the prices deserve a word. The resource rents die: $p^{S*}=p^{X*}=0$, because the dividends both stocks pay are proportional to activity that has ceased. And $p^{W*}$ need not vanish even though the stockpile empties: the per-tonne margins remain well defined in the limit, and at a rest point the asset price of a stockpiled tonne is its handling dividend, discounted for the geometric wait until handling, $p^{\mathcal W*}=\mu h^*/\left(\rho+\mu\right)$, which is exactly the rest-point waste margin~\eqref{eq:sp:sum:W-limit}. The workhorse solves the rest-point block $\left(p^{W*},x^*,\varpi^*\right)$ explicitly, including the corner in which the last waste is treated only for the diversion motive --- and, with the collection charge on the treated tonne, the further corner in which even that is not worth paying for and handling ends untreated.

\smalltitle{How the material era ends ($\bar R=0$).} The marginal notion of essentiality splits case (i) into two economically distinct endings, and governs the terminal behaviour of the recycling block in all cases.

\emph{Ending by choice (finite choke value).} If $F^\infty_R\left(K^*,0,0\right)<\infty$, the value of a marginal tonne at the rest point is finite, and the extraction corner~\eqref{eq:sp:sum:N-corner} can hold with a positive reserve. Extraction ceases in finite time, when the choke condition first binds, and reserves are \emph{stranded}: $S_\infty>0$ generically. The material era ends the way the stone age is said to have ended --- not for lack of stones. The recycling margins wind down with it: the virgin displacement value $V$ stays bounded, and if the marginal product of the first unit of recycling capital falls short, recycling hits its corner and treatment survives, if at all, on the diversion motive alone.

\emph{Ending by exhaustion (infinite marginal value).} If instead $F^\infty_R\to\infty$ as $R\to0$, no finite choke exists: extraction continues forever at vanishing rates, the interior condition~\eqref{eq:sp:sum:N} requires marginal extraction cost to diverge with $\Psi$, and reserves are squeezed to zero, $S\to0$. The same divergence drives the recycling block to its technological limit: $V\to\infty$ forces $x\to\infty$, the yield climbs to its ceiling, and the treatment share is pushed into its upper branch. \emph{Maximal circularity is the shadow of maximal scarcity}: the economy becomes as circular as technology allows precisely because the alternative to a recycled tonne is becoming infinitely expensive. Consumption still converges to the same $C^*$ of~\eqref{eq:sp:lr:mgr} in the level-inessential case --- the squeeze concerns the vanishing material margin, not the goods block.

Case (ii) inherits the exhaustion ending with nothing to converge to: the classical exhaustible-resource collapse, with two twists --- the ceiling $\bar X_{\max}$ makes exploration a temporary reprieve rather than an escape, and recycling multiplies the budget by at most $\left(1-\bar a\right)^{-1}$ without changing the verdict.

\subsection{The floorless benchmark: trending technology}\label{sec:sp:longrun:trend}

With $\bar R=0$, sustained progress overturns the dichotomy: survival no longer depends on essentiality at all. The material budget constrains cumulated \emph{physical} input, and sustained progress converts a shrinking physical flow into non-shrinking production --- either because effective materials $B_tR$ hold up as $R$ decays ($g_B>0$), or because TFP growth outpaces the drag from a declining material input ($g_A>0$).

\begin{definition}[balanced dematerialization path]\label{def:lr:bdp}
A path is a balanced dematerialization path (BDP, state B2) if, asymptotically, the goods block grows at a constant rate $g\ge0$ --- $C,K,Y$ at rate $g$, $\ln\left(1+r\right)\to\ln\left(1+\rho\right)+\eta g$ --- while the material block decays exponentially: $R,N,W,H,T,M^K$ at rate $\nu>0$, with $S$ and $\bar X_{\max}-X$ converging to zero, and the recycling margins converging: $a\to\bar a$, $\varpi$ into its upper branch.
\end{definition}

On a BDP every constraint of Section~\ref{sec:sp:setup:primitives} is respected with room to spare: exponentially decaying flows are summable, so the cumulative bounds bind only in the limit, and $\Xi\to0$ with $P\to0$. The waste stock rides along rather than binding: with $\mu>1-e^{-\nu}$, $\mathcal W_t=W_t/\left[\mu-\left(1-e^{-\nu}\right)\right]\cdot\left(1+o\left(1\right)\right)$ decays at the same rate $\nu$ as the flows; if instead $\mu\le1-e^{-\nu}$ the stockpile is the slowest-decaying material reservoir and its drainage --- geometric at factor $1-\mu$, rate $-\ln\left(1-\mu\right)$ --- sets the terminal decay of the whole material block. The asymptotic optimality conditions have the structure of the exhaustion ending stretched out forever: the material value $\Psi$ grows at rate $g+\nu$, matched by diverging marginal extraction and discovery costs; the recycling-capital condition forces $a\to\bar a$ and full treatment follows from $c^{T\prime}\left(1\right)<\infty$.

Two features of the price system are worth recording, and both are consequences of the yield ceiling. First, the waste price turns negative and grows in magnitude with $\Psi$: from the rest-point form of the waste margin, $p^W\approx-\bar a\varpi\,\Psi/\left(1-\bar a\varpi\right)<0$ asymptotically, up to a bounded correction factor that the balanced-growth recursion determines (the workhorse records it) --- the stockpile is an asset, a store of future material valued at the multiplied worth of what it contains. Second, despite $\zeta$ growing with $\left|p^W\right|$, every wedge $\zeta\Omega^j$ remains bounded, because $\Omega\propto R/Y$ decays at the rate material prices grow. Under a hard ceiling the accounting block therefore stays regular along the entire path; under the soft ceiling the same expression makes the multiplier outrun $\Psi$ as $a\varpi\to1$, which is why the soft-ceiling alternative to a BDP is not a faster BDP but the closed loop of Section~\ref{sec:sp:longrun:floor}. With $\bar R=0$ the planner genuinely faces that choice, and~\eqref{eq:sp:lr:bdp-vs-c} below settles it.

The growth rate on a BDP is determined jointly with the decay rate, by the production side and the extraction margin together; in the workhorse's Cobb--Douglas case the two rows are $\left(1-\beta_K\right)g=g_A+\gamma\left(g_B-\nu\right)$ and $\left(\mu_N-1\right)\nu=g$, giving
\begin{align}
g &= \frac{g_A+\gamma\left(g_B-\nu\right)}{1-\beta_K}
\;=\;\frac{\left(\mu_N-1\right)\left(g_A+\gamma g_B\right)}{\left(\mu_N-1\right)\left(1-\beta_K\right)+\gamma},
&
\nu &= \frac{g_A+\gamma g_B}{\left(\mu_N-1\right)\left(1-\beta_K\right)+\gamma} .
\label{eq:sp:lr:bdp-growth}
\end{align}
The stock-effect elasticity of extraction costs, $\mu_N$, is what selects the regime: $\mu_N>1$ gives sustained growth, $\mu_N=1$ gives a sustained consumption \emph{level} --- material-augmenting progress exactly offsetting physical decay with effective abundance --- and $\mu_N<1$ gives managed decline, with the exponential family failing altogether once $\gamma\le\left(1-\mu_N\right)\left(1-\beta_K\right)$, where extraction ends in finite time and the path continues on recycled material alone. Appendix~\ref{sec:app:workhorse:bdp} carries the full rate table and the derivation.

Comparing~\eqref{eq:sp:lr:bdp-growth} with the circular rate~\eqref{eq:sp:lr:c2-growth} settles the selection question, since the two share a numerator:
\begin{align}
g_{\mathrm{B2}} &= \frac{\left(\mu_N-1\right)\left(1-\beta_K\right)}{\left(\mu_N-1\right)\left(1-\beta_K\right)+\gamma}\;g_{\mathrm{C}}\;<\;g_{\mathrm{C}} .
\label{eq:sp:lr:bdp-vs-c}
\end{align}
\emph{Balanced dematerialization always grows strictly more slowly than the closed loop}, and the entire gap is the physical decay $\nu$ that the squeeze path accepts and the loop does not. Since closing the loop costs an asymptotically vanishing share of the capital stock, the planner closes it whenever the yield tail allows: state C dominates B2 wherever both are feasible, and the floorless soft-ceiling cell of the taxonomy is a C cell, not a comparison.

Finally, the floor is what closes this door: with $\bar R>0$ no BDP exists, because $R$ cannot decay through the floor while production continues --- which is exactly why the floor cells of Section~\ref{sec:sp:longrun:floor} offer only the stark choice between the closed loop and collapse.

\subsection{Taxonomy}\label{sec:sp:longrun:taxonomy}

The classification is summarised below. ``Ceiling'' distinguishes the hard ceiling (H) from the soft ceiling (S), with which we group only those soft ceilings that fail the closure criterion of Appendix~\ref{sec:app:workhorse:cbgp} --- tails so slow that the yield elasticity $\epsilon_a$ vanishes, of which the logarithmic family is the leading example. The exponential and power families both close the loop, so the column is a genuine ceiling distinction and no longer a conjecture about tails. ``Circularity'' is the asymptotic recycled share $a\varpi$. Pollution converges to zero in every row; in the collapse rows with $\bar R>0$ it does so after the legacy emission stream of the draining stockpile.

\begin{center}
\begin{tabular}{llllll}
\toprule
$\bar R$ & Ceiling & Regime & Essentiality (level) & State & Long-run $C$ \\
\midrule
$0$ & hard & (S) & inessential & B1: dematerialized rest point & $C^*>0$, eq.~\eqref{eq:sp:lr:mgr} \\
$0$ & hard & (S) & essential & B3: exhaustion collapse & $C\to0$ \\
$0$ & hard & (T) & any & B2: balanced dematerialization & grows at $g\ge0$, eq.~\eqref{eq:sp:lr:bdp-growth} \\
$0$ & soft & (S) & either & B1 or B3 (loop blocked by Prop.~\ref{prop:lr:growth}) & as above \\
$0$ & soft & (T) & any & C, dominating B2 by~\eqref{eq:app:wh:bdp-vs-c} & grows at $g>0$ \\
\midrule
$>0$ & hard & (S) & --- & A: shutdown collapse & $C\to0$ \\
$>0$ & hard & (T) & --- & A: collapse \emph{despite growth} & $C\to0$ \\
$>0$ & soft & (S) & --- & A (loop blocked by Prop.~\ref{prop:lr:growth}) & $C\to0$ \\
$>0$ & soft & (T) & --- & C: perpetual circular growth, if $\mathcal M_\infty>\bar R\mathcal T$ & grows at $g>0$ \\
\bottomrule
\end{tabular}
\end{center}

Three readings organise the table.

\emph{First, the floor deletes the essentiality column.} With $\bar R>0$, level essentiality holds by construction and marginal essentiality is pre-empted by the discrete shutdown, which stops material use at a strictly positive flow, so the exhaustion-versus-choke distinction that organises the floorless benchmark decides nothing; the classification depends only on the ceiling and the regime. This is a large simplification, and it is the floor doing it.

\emph{Second, growth is not a free pass.} Trending technology rescues the floorless economy unconditionally (row B2) but rescues the floor economy only through the closed loop: under a hard ceiling the material era has bounded duration whatever $g_A$ and $g_B$ are, because progress augments effective materials and cannot touch the floor. The cell (H)$\times$(T)$\times\bar R>0$ --- a technologically brilliant economy collapsing on schedule --- is the taxonomy's starkest entry.

\emph{Third, the bottom-right cell is the paper's sharpest claim.} With a floor, survival and perpetual circularity are the same thing, available only under the soft ceiling, affordable only because the economy grows (Proposition~\ref{prop:lr:growth}), and reachable only if the economy retains enough matter to clear the floor at its own turnover rate, $\mathcal M_\infty>\bar R\mathcal T$ by~\eqref{eq:sp:lr:little}. In every other row the circular economy is a transition technology: it determines how much value the budget $\mathcal B$ yields and how much pollution damage is suffered en route --- the quantitative Part of the paper measures exactly that --- while in this one cell it is the difference between growth forever and collapse. Meanwhile the scarcity reading of circularity survives throughout: full treatment and ceiling-level yields emerge asymptotically wherever materials stay scarce, and precisely because they stay scarce; the one economy that abandons circularity is the one that outgrows materials at a finite choke value.

\emph{A fourth reading, on what the classification is not.} Two columns of the table are primitives that no policy moves --- the floor and the ceiling --- and one, the regime, is a statement about trends. Nothing in the table is a choice. What \emph{is} a choice is which side of the survival inequality~\eqref{eq:sp:lr:little} an economy in the bottom-right cell ends up on, since $\mathcal M_\infty$ is the endogenous outcome of every extraction and leakage decision made along the way. The taxonomy classifies the destinations; the transition decides which one is reached, and that is a quantitative question.

\smalltitle{Configurations to check, not to classify.} Three path properties cut across the table. \emph{Skiba points:} the pollution non-convexity ($\theta'<0$) can strand the economy at an inferior rest point; we condition on the limit attained. \emph{Tipping:} $r+\theta+\theta'P\le0$ makes $p^P$ diverge; excluded by Assumption~\ref{ass:lr:theta}, which is not innocuous. \emph{Utility degeneracy:} the collapse rows cannot be ranked internally when $u\left(0\right)=-\infty$; see the scope remark in Section~\ref{sec:sp:longrun:scope}. Appendix~\ref{sec:app:workhorse:shutdown} removes part of the last difficulty by showing that consumption after a shutdown declines geometrically rather than vanishing, so collapse paths have a finite value whenever $\beta\gamma_C^{1-\eta}<1$.

\smalltitle{Open items.} In decreasing order of importance. (i)~Sufficiency, which Appendix~\ref{sec:app:suff} takes as far as it goes. Given the operating set and the material-intensity path, the necessary conditions above are sufficient under four checkable conditions on primitives, and transversality reduces to a single condition on capital because mass conservation bounds the other five states. What is not established is optimality against deviations that change the intensity path, optimality across operating sets where a shutdown might not be permanent, and sufficiency where extraction has a finite choke value --- which is exactly the configuration in which stranding, and hence multiplicity, is possible. (ii)~The transition into state C, in particular how much of the budget an optimal path retains --- which by~\eqref{eq:sp:lr:little} fixes the level of the entire future, and which only the quantitative model can answer. (iii)~The regularity restrictions the state-C price block imposes on calibrations (Appendix~\ref{sec:app:workhorse:cbgp}), which are inequalities on wedges rather than conditions on primitives. (iv)~The squeeze-case decay rates in regime (S) with $1<\varsigma<\infty$, and the rate algebra away from Cobb--Douglas. (v)~The irreversible-pollution variant, which adds a hysteresis dimension to every row.
